\documentclass[12pt, aspectratio=169]{beamer}

\input{../header}
\usepackage{booktabs}

\title{Regression Analysis}
\author{Md. Aminul Islam Shazid}
\date{}


\begin{document}
    {
        \setbeamertemplate{footline}{}
        \addtocounter{framenumber}{-2}

        \begin{frame}
            \titlepage
        \end{frame}

        \begin{frame}{Outline}
            \frametitle<1>{Table of Contents}
            \frametitle<2>{Table of Contents (cont.)}
            \vfill
            % \tableofcontents[subsectionstyle=hide]
            \only<1>{\tableofcontents[sections={1-4},hideallsubsections]}
            \only<2>{\tableofcontents[sections={5-8},hideallsubsections]}
            \vfill
        \end{frame}
    }

    %------------------------------------------------------------
    \section{Introduction}
    %------------------------------------------------------------

    \begin{frame}{Regression Analysis --- Overview}
        \begin{itemize}
            \item Regression analysis examines the \textbf{statistical relationship}
                between two or more variables
            \item One variable is designated the \textbf{dependent variable} ($Y$,
                the variable to be explained or predicted), and the remaining
                variables are designated \textbf{independent variables}
                ($X_1, X_2, \ldots$, the explanatory variables)
            \item Regression \textbf{models} and \textbf{predicts} the average
                value of $Y$ for given values of $X$
            \item \textbf{Important:} regression describes a statistical
                association, not necessarily a causal one --- the designation of
                a variable as ``independent'' reflects the analyst's choice of
                model, not a proven direction of causation
        \end{itemize}
    \end{frame}


    \begin{frame}{Examples}
        \begin{itemize}
            \item Number of family members ($X$) and monthly food expenditure ($Y$, thousand taka)
            \item Weather temperature as \(X\) and sales of ice cream \(Y\)
            \item Income as \(X\) and expenditure as \(Y\)
        \end{itemize}
    \end{frame}


    %------------------------------------------------------------
    \section{Simple Linear Regression}
    %------------------------------------------------------------

    \begin{frame}{Simple Linear Regression Model}
        \textbf{Population model:}
        \[
            Y_i = \alpha + \beta X_i + \epsilon_i, \quad i = 1, 2, \ldots, n
        \]
        \begin{itemize}
            \item $Y$ --- dependent variable
            \item $X$ --- independent variable
            \item $\alpha$ --- intercept (regression parameter)
            \item $\beta$ --- slope (regression parameter)
            \item $\epsilon_i$ --- error term (unexplained variation)
        \end{itemize}
    \end{frame}


    \begin{frame}{Estimated Model}
        \begin{itemize}
            \item The population model is estimated using data from sample
            \item Sample (estimated) model:
            \[
                \hat{y}_i = a + b\,x_i
            \]
            so that the \textbf{residual} $e_i = y_i - \hat{y}_i$
            \item \(\hat{y}_i\) are the predicted values, \(a\) is the estimate for \(\alpha\) and \(b\) is the estimate for \(\beta\)
        \end{itemize}
    \end{frame}


    \begin{frame}{Assumptions of Simple Linear Regression}
        \begin{enumerate}
            \item $X$ values are fixed (non-random)
            \item The relationship between $X$ and $Y$ is \textbf{linear}
            \item Error terms are normally distributed with mean zero and constant variance:
                \[
                    \epsilon_i \sim N(0,\,\sigma^2)
                \]
            \item $X$ and $\epsilon$ are uncorrelated:
                $\mathrm{Corr}(X, \epsilon) = 0$
        \end{enumerate}
    \end{frame}


    \begin{frame}{Estimation --- Least Squares Method}
        \begin{itemize}
            \item \textbf{Principle:} find $a$ and $b$ that minimise the
                \emph{sum of squared errors} (SSE):
                \[
                    \min_{a,b}\;\sum_{i=1}^{n} e_i^2
                    = \min_{a,b}\;\sum_{i=1}^{n}(y_i - \hat{y}_i)^2
                \]
            \item \textbf{Least squares estimates (LSE):}
                \begin{align*}
                    b &= \frac{\sum (x_i - \bar{x})(y_i - \bar{y})}{\sum (x_i - \bar{x})^2} = \frac{n\sum xy - \sum x \sum y}{n\sum x^2 - (\sum x)^2} \\[10pt]
                    a &= \frac{\sum y}{n} - b\,\frac{\sum x}{n} = \bar{y} - b\bar{x}
                \end{align*}
        \end{itemize}
    \end{frame}


    \begin{frame}{SLR Derivation --- The Least Squares Criterion}
        Our goal is to find $a$ and $b$ that minimize the \textbf{Sum of Squared Errors (SSE)}:
        \begin{equation*}
            S(a, b) = \sum_{i=1}^{n} e_i^2 = \sum_{i=1}^{n} (y_i - \hat{y}_i)^2 = \sum_{i=1}^{n} (y_i - a - bx_i)^2
        \end{equation*}

        To find the minimum, we take partial derivatives with respect to $a$ and $b$ and set them to zero:
        \begin{align}
            \frac{\partial S}{\partial a} &= -2 \sum_{i=1}^{n} (y_i - a - bx_i) = 0 \notag \\
            \frac{\partial S}{\partial b} &= -2 \sum_{i=1}^{n} x_i (y_i - a - bx_i) = 0 \label{eq:deriv2}
        \end{align}
    \end{frame}


    \begin{frame}{SLR Derivation --- Solving the Normal Equations}
    From the first derivative, we get the first \textbf{Normal Equation}:
    \begin{equation*}
        \sum y_i = na + b \sum x_i \implies \bar{y} = a + b\bar{x} \implies a = \bar{y} - b\bar{x}
    \end{equation*}

        Substitute $a = \bar{y} - b\bar{x}$ into equation \ref{eq:deriv2}:
        \begin{align*}
            & \sum x_i (y_i - (\bar{y} - b\bar{x}) - bx_i) = 0 \\
            \rightarrow & \sum x_i (y_i - \bar{y}) - b \sum x_i (x_i - \bar{x}) = 0
        \end{align*}
        Solving for $b$ gives the standard estimator:
        \begin{equation*}
            b = \frac{\sum (x_i - \bar{x})(y_i - \bar{y})}{\sum (x_i - \bar{x})^2} = \frac{n\sum xy - \sum x \sum y}{n\sum x^2 - (\sum x)^2}
        \end{equation*}
    \end{frame}


    \begin{frame}{Interpretation of Estimated Parameters}
        \begin{columns}
            \begin{column}{0.55\textwidth}
                \begin{itemize}
                    \item \textbf{Intercept} $a$: the estimated value of $Y$ when $X = 0$ (baseline level)
                    \item \textbf{Slope} $b$: for every \textbf{1-unit increase} in $X$, the average change in $Y$ is $b$ units
                        \begin{itemize}
                            \item $b > 0$: $Y$ increases on average
                            \item $b < 0$: $Y$ decreases on average
                        \end{itemize}
                \end{itemize}
            \end{column}
            \begin{column}{0.45\textwidth}
                \begin{center}
                    % Schematic diagram of the regression line
                    \begin{tikzpicture}[scale=0.85]
                        \draw[->] (0,0) -- (4.2,0) node[right] {$x$};
                        \draw[->] (0,0) -- (0,3.5) node[above] {$y$};
                        \draw[thick, blue] (0,0.6) -- (4,3.4)
                            node[right] {$\hat{y}=a+bx$};
                        \draw[dashed] (0,0.6) -- (-0.1,0.6)
                            node[left] {$a$};
                        % slope triangle
                        \draw[<->] (1.5,1.35) -- (2.5,1.35)
                            node[midway, below] {$1$};
                        \draw[<->] (2.5,1.35) -- (2.5,2.05)
                            node[midway, right] {$b$};
                    \end{tikzpicture}
                \end{center}
            \end{column}
        \end{columns}
    \end{frame}


    %------------------------------------------------------------
    \section{Goodness of Fit}
    %------------------------------------------------------------

    \begin{frame}{Goodness of Fit --- $R^2$}
        \textbf{Decomposition of total variation:}
        \[
            \underbrace{\sum(y_i - \bar{y})^2}_{\text{TSS (Total)}}
            = \underbrace{\sum(\hat{y}_i - \bar{y})^2}_{\text{SSR (Regression)}}
            + \underbrace{\sum(y_i - \hat{y}_i)^2}_{\text{SSE (Error)}}
        \]

        \vspace{0.5em}

        \textbf{Coefficient of determination:}
        \[
            R^2 = \frac{\text{SSR}}{\text{TSS}} = 1 - \frac{\text{SSE}}{\text{TSS}}
            = 1 - \frac{\sum e_i^2}{\sum y_i^2 - \dfrac{(\sum y_i)^2}{n}}
        \]

        \textbf{Interpretation:}
        \begin{itemize}
            \item $0 \leqslant R^2 \leqslant 1$
            \item $R^2 \to 0$: poor fit; $R^2 \to 1$: good fit
            \item ``$R^2 \times 100$\,\% of the variation in $Y$ is explained
                by the variation in $X$''
        \end{itemize}
    \end{frame}


    %------------------------------------------------------------
    \section{Example (Simple Linear Regression)}
    %------------------------------------------------------------

    \begin{frame}{Example --- Family Members vs.\ Food Expenditure}
        \textbf{Problem:} Fit a regression of monthly food expenditure ($y$, thousand
        taka) on number of family members ($x$). Interpret the estimates and
        compute $R^2$.  Predict expenditure when $x = 10$.

        \begin{center}
            \small
            \begin{tabular}{ccccc}
                \toprule
                $x$ & $y$ & $x^2$ & $xy$ & $y^2$ \\
                \midrule
                2 & 5  & 4  & 10 & 25  \\
                3 & 7  & 9  & 21 & 49  \\
                6 & 11 & 36 & 66 & 121 \\
                4 & 8  & 16 & 32 & 64  \\
                7 & 13 & 49 & 91 & 169 \\
                3 & 6  & 9  & 18 & 36  \\
                6 & 12 & 36 & 72 & 144 \\
                \midrule
                $\sum x = 31$ & $\sum y = 62$ & $\sum x^2 = 159$ &
                $\sum xy = 310$ & $\sum y^2 = 608$ \\
                \bottomrule
            \end{tabular}
        \end{center}
    \end{frame}


    \begin{frame}{Example --- Solution}
        \textbf{Least squares estimates} ($n = 7$):
        \[
            b = \frac{7(310) - 31(62)}{7(159) - 31^2} = 1.63,
            \qquad
            a = \frac{62}{7} - 1.63 \times \frac{31}{7} = 1.64
        \]

        \textbf{Estimated regression equation:}
        \[
            \hat{y} = 1.64 + 1.63\,x
        \]

        \textbf{Interpretation:}
        \begin{itemize}
            \item $a = 1.64$: when no family members ($x=0$), baseline food
                expenditure is 1.64 thousand taka
            \item $b = 1.63$: adding one family member increases average food
                expenditure by 1.63 thousand taka
        \end{itemize}
    \end{frame}


    \begin{frame}{Example --- $R^2$ and Prediction}
        \textbf{Residuals:} $\sum e_i^2 = 1.05$

        \[
            R^2 = 1 - \frac{1.05}{608 - \dfrac{62^2}{7}} = 0.9821
        \]

        \textbf{Interpretation:} 98.21\% of the variation in monthly food
        expenditure is explained by the number of family members.
        The model has an excellent fit.

        \vspace{0.8em}

        \textbf{Prediction} (for $x = 10$ family members):
        \[
            \hat{y} = 1.64 + 1.63 \times 10 = 17.93 \text{ thousand taka}
        \]
    \end{frame}


    %------------------------------------------------------------
    \section{Test of Hypothesis for Regression Parameters}
    %------------------------------------------------------------

    \begin{frame}{Why Test Regression Parameters?}
        \begin{itemize}
            \item The slope $b$ is estimated from a \emph{sample}; we need to
                know whether a non-zero $b$ reflects a true population
                relationship or just random sampling variation
            \item \textbf{Testing $H_0: \beta = 0$} asks:
                ``Does $X$ have a statistically significant linear effect on $Y$?''
            \item If we fail to reject $H_0: \beta = 0$, the regression line
                adds no predictive value over simply using $\bar{y}$
            \item Two complementary approaches:
            \begin{itemize}
                \item \textbf{$t$-test} for individual coefficients
                \item \textbf{$F$-test} for overall significance of the model
            \end{itemize}
        \end{itemize}
    \end{frame}


    \begin{frame}{Regression ANOVA Table}
        \begin{itemize}
            \item Just as in one-way ANOVA, we partition the total variation:
                \[
                    \text{TSS} = \text{SSR} + \text{SSE}
                \]
            \item Mean squares:
                \[
                    \text{MSR} = \frac{\text{SSR}}{1},
                    \qquad
                    \text{MSE} = \frac{\text{SSE}}{n-2}
                    \quad \text{(estimated } \sigma^2 \text{)}
                \]
        \end{itemize}

        \vspace{0.4em}

        \begin{center}
            \begin{tabular}{lcccc}
                \toprule
                \textbf{Source} & \textbf{df} & \textbf{SS} & \textbf{MS} & \textbf{F} \\
                \midrule
                Regression & $1$   & SSR & MSR & $\text{MSR}/\text{MSE}$ \\
                Error      & $n-2$ & SSE & MSE & \\
                Total      & $n-1$ & TSS & & \\
                \bottomrule
            \end{tabular}
        \end{center}
    \end{frame}


    \begin{frame}{$F$-test for Overall Significance}
        \begin{itemize}
            \item \textbf{Hypotheses:}
                \[
                    H_0: \beta = 0 \quad \text{vs} \quad H_A: \beta \neq 0
                \]
            \item \textbf{Test statistic:}
                \[
                    F_{\text{cal}} = \frac{\text{MSR}}{\text{MSE}}
                    \sim F_{1,\,n-2} \quad \text{under } H_0
                \]
            \item \textbf{Decision rule:}
                \begin{itemize}
                    \item Reject $H_0$ if $F_{\text{cal}} > F_{\alpha,\,1,\,n-2}$
                    \item Or equivalently, reject if $p\text{-value} \leqslant \alpha$
                \end{itemize}
            \item \textbf{Note:} For simple linear regression (one predictor),
                the $F$-test and the two-tailed $t$-test for $\beta$ are
                equivalent; $F_{\text{cal}} = t_{\text{cal}}^2$
        \end{itemize}
    \end{frame}


    \begin{frame}{$t$-test for the Slope $\beta$}
        \begin{itemize}
            \item \textbf{Hypotheses:} \quad $H_0: \beta = 0$ \; vs \;
                $H_A: \beta \neq 0$ \quad (two-tailed)
            \item \textbf{Standard error of $b$:}
                \[
                    \mathrm{SE}(b) = \sqrt{\frac{\text{MSE}}{\sum(x_i - \bar{x})^2}}
                    = \frac{s}{\sqrt{\sum(x_i-\bar{x})^2}},
                    \quad s = \sqrt{\text{MSE}}
                \]
            \item \textbf{Test statistic:}
                \[
                    t_{\text{cal}} = \frac{b - 0}{\mathrm{SE}(b)} \sim t_{n-2}
                    \quad \text{under } H_0
                \]
            \item \textbf{Decision:} reject $H_0$ if $|t_{\text{cal}}| > t_{\alpha/2,\,n-2}$
        \end{itemize}
    \end{frame}


    \begin{frame}{$t$-test for the Intercept $\alpha$}
        \begin{itemize}
            \item \textbf{Hypotheses:} \quad $H_0: \alpha = 0$ \; vs \;
                $H_A: \alpha \neq 0$
            \item \textbf{Standard error of $a$:}
                \[
                    \mathrm{SE}(a) = s\sqrt{\frac{1}{n} + \frac{\bar{x}^2}{\sum(x_i - \bar{x})^2}}
                \]
            \item \textbf{Test statistic:}
                \[
                    t_{\text{cal}} = \frac{a - 0}{\mathrm{SE}(a)} \sim t_{n-2}
                    \quad \text{under } H_0
                \]
            \item In practice, the intercept test is often less relevant than the slope test
        \end{itemize}
    \end{frame}


    \begin{frame}{Example --- Hypothesis Test for the Slope}
        Using the family-expenditure data ($n = 7$, $b = 1.63$, $\text{SSE} = 1.05$):

        \[
            \text{MSE} = \frac{1.05}{7-2} = 0.21,
            \qquad
            s = \sqrt{0.21} = 0.458
        \]

        \[
            \sum(x_i - \bar{x})^2 = \sum x_i^2 - \frac{(\sum x_i)^2}{n}
            = 159 - \frac{31^2}{7} = 22.86
        \]

        \[
            \mathrm{SE}(b) = \frac{0.458}{\sqrt{22.86}} = 0.096,
            \qquad
            t_{\text{cal}} = \frac{1.63}{0.096} = 16.98
        \]

        \textbf{Critical value:} $t_{0.025,\,5} = 2.571$

        \vspace{0.25em}

        \textbf{Decision:} Since $|t_{\text{cal}}| = 16.98 > 2.571$, we
        \textbf{reject} $H_0: \beta = 0$ at 5\% level of significance.
        There is a statistically significant linear relationship between the number
        of family members and monthly food expenditure.
    \end{frame}


    \begin{frame}{Confidence Interval for $\beta$}
        \begin{itemize}
            \item A $100(1-\alpha)\%$ confidence interval for the true slope $\beta$:
                \[
                    b \pm t_{\alpha/2,\,n-2}\;\mathrm{SE}(b)
                \]
            \item \textbf{Example} (95\% CI, family-expenditure data):
                \[
                    1.63 \pm 2.571 \times 0.096 = 1.63 \pm 0.247
                    = (1.383,\;1.877)
                \]
            \item \textbf{Interpretation:} we are 95\% confident that for each
                additional family member, mean monthly food expenditure increases
                by between 1.383 and 1.877 thousand taka
            \item If the interval does not contain 0, this is equivalent to
                rejecting $H_0: \beta = 0$ at the 5\% level
        \end{itemize}
    \end{frame}


    %------------------------------------------------------------
    \section{Multiple Linear Regression}
    %------------------------------------------------------------

    \begin{frame}{Multiple Linear Regression --- Model}
        When there are $k$ independent variables $X_1, X_2, \ldots, X_k$:
        \[
            Y_i = \beta_0 + \beta_1 X_{1i} + \beta_2 X_{2i}
            + \cdots + \beta_k X_{ki} + \epsilon_i
        \]

        \begin{itemize}
            \item $\beta_0$ --- intercept
            \item $\beta_j$ --- \textbf{partial regression coefficient} of $X_j$;
                the expected change in $Y$ per unit increase in $X_j$,
                holding all other $X$'s constant
            \item \textbf{Assumption:} no strong linear correlation (multicollinearity)
                among the $X$ variables
        \end{itemize}

        \textbf{Matrix form:}
        \[
            \mathbf{Y} = \mathbf{X}\boldsymbol{\beta} + \boldsymbol{\epsilon}
            \qquad \Rightarrow \qquad
            \hat{\boldsymbol{\beta}} = (\mathbf{X}'\mathbf{X})^{-1}\mathbf{X}'\mathbf{y}
        \]
    \end{frame}


    \begin{frame}{Matrix Notation --- What Each Object Represents}
        \begin{columns}
            \begin{column}{0.48\textwidth}
                $\mathbf{Y}$ is an $n \times 1$ \textbf{column vector} of observed
                responses:
                \vspace{-0.5em}
                \[
                    \mathbf{Y} =
                    \begin{pmatrix} Y_1 \\ Y_2 \\ \vdots \\ Y_n \end{pmatrix}_{n \times 1}
                \]
                $\boldsymbol{\beta}$ is a $(k+1) \times 1$ vector of unknown parameters:
                \vspace{-0.5em}
                \[
                    \boldsymbol{\beta} =
                    \begin{pmatrix} \beta_0 \\ \beta_1 \\ \vdots \\ \beta_k \end{pmatrix}_{(k+1)\times 1}
                \]
            \end{column}
            \begin{column}{0.52\textwidth}
                $\mathbf{X}$ is the $n \times (k+1)$ \textbf{design matrix}. The
                \textbf{first column is all 1s} (for the intercept $\beta_0$);
                the remaining $k$ columns contain the observed values of
                $X_1, \ldots, X_k$:
                \[
                    \mathbf{X} =
                    \begin{pmatrix}
                        1 & x_{11} & x_{21} & \cdots & x_{k1} \\
                        1 & x_{12} & x_{22} & \cdots & x_{k2} \\
                        \vdots & \vdots & \vdots & \ddots & \vdots \\
                        1 & x_{1n} & x_{2n} & \cdots & x_{kn}
                    \end{pmatrix}_{n\times(k+1)}
                \]

                $\boldsymbol{\epsilon}$ is an $n \times 1$ vector of errors,
                $\boldsymbol{\epsilon} \sim N(\mathbf{0},\,\sigma^2 \mathbf{I})$
            \end{column}
        \end{columns}
    \end{frame}


    \begin{frame}{MLR Derivation --- The Matrix Approach}
        In matrix form, the model is $\mathbf{y} = \mathbf{X}\boldsymbol{\beta} + \boldsymbol{\epsilon}$. We minimize the residual sum of squares:
        \begin{equation*}
            S(\boldsymbol{\beta}) = \boldsymbol{\epsilon}'\boldsymbol{\epsilon} = (\mathbf{y} - \mathbf{X}\boldsymbol{\beta})'(\mathbf{y} - \mathbf{X}\boldsymbol{\beta})
        \end{equation*}

        \begin{block}{Expansion and Gradient}
            Expanding the quadratic form:
            \begin{equation*}
                S(\boldsymbol{\beta}) = \mathbf{y}'\mathbf{y} - 2\boldsymbol{\beta}'\mathbf{X}'\mathbf{y} + \boldsymbol{\beta}'\mathbf{X}'\mathbf{X}\boldsymbol{\beta}
            \end{equation*}
            Taking the gradient with respect to $\boldsymbol{\beta}$ and setting to $\mathbf{0}$:
            \begin{equation*}
                \frac{\partial S}{\partial \boldsymbol{\beta}} = -2\mathbf{X}'\mathbf{y} + 2\mathbf{X}'\mathbf{X}\boldsymbol{\beta} = \mathbf{0}
            \end{equation*}
        \end{block}
    \end{frame}


    \begin{frame}{MLR Derivation --- The Matrix Approach (cont.)}
            Rearranging gives the \textbf{Normal Equations}: $\mathbf{X}'\mathbf{X}\hat{\boldsymbol{\beta}} = \mathbf{X}'\mathbf{y}$.

            \vspace{1em}

            Assuming $\mathbf{X}'\mathbf{X}$ is non-singular: \quad $\mathbf{\hat{\boldsymbol{\beta}} = (X'X)^{-1}X'y}$
    \end{frame}


    \begin{frame}{Multiple Regression --- $F$-test for Overall Significance}
        \begin{itemize}
            \item \textbf{Hypotheses:}
                \[
                    H_0: \beta_1 = \beta_2 = \cdots = \beta_k = 0
                    \quad \text{vs} \quad
                    H_A: \text{at least one } \beta_j \neq 0
                \]
            \item ANOVA table for multiple regression:
        \end{itemize}
        \begin{center}
            \begin{tabular}{lcccc}
                \toprule
                \textbf{Source} & \textbf{df} & \textbf{SS} & \textbf{MS} & \textbf{F} \\
                \midrule
                Regression & $k$     & SSR & $\text{SSR}/k$     & $\text{MSR}/\text{MSE}$ \\
                Error      & $n-k-1$ & SSE & $\text{SSE}/(n-k-1)$ & \\
                Total      & $n-1$   & TSS & & \\
                \bottomrule
            \end{tabular}
        \end{center}
        \[
            F_{\text{cal}} = \frac{\text{MSR}}{\text{MSE}} \sim F_{k,\,n-k-1}
            \quad \text{under } H_0
        \]
        \textbf{Decision:} reject $H_0$ if $F_{\text{cal}} > F_{\alpha,\,k,\,n-k-1}$
    \end{frame}


    \begin{frame}{Multiple Regression --- $t$-tests for Individual Coefficients}
        \begin{itemize}
            \item For each $\beta_j$ ($j = 1, 2, \ldots, k$):
                \[
                    H_0: \beta_j = 0 \quad \text{vs} \quad H_A: \beta_j \neq 0
                \]
            \item \textbf{Test statistic:}
                \[
                    t_{\text{cal}} = \frac{b_j}{\mathrm{SE}(b_j)} \sim t_{n-k-1}
                    \quad \text{under } H_0
                \]
            \item \textbf{Decision:} reject $H_0$ if
                $|t_{\text{cal}}| > t_{\alpha/2,\,n-k-1}$
            \item A non-significant $t$-test for $\beta_j$ suggests that
                $X_j$ may not contribute significantly to $Y$ given all other
                predictors are in the model
            \item These individual tests should be interpreted \emph{after}
                the overall $F$-test confirms at least one predictor is useful
        \end{itemize}
    \end{frame}


    \begin{frame}{Multiple Regression --- Example (Automobile Data)}
        \textbf{Model:} Mileage ($y$) on displacement ($x_1$), horsepower ($x_2$),
        weight ($x_{10}$), and transmission type (\texttt{Auto})

        \[
            \hat{y} = 34.5 - 0.05\,x_1 + 0.03\,x_2 - 0.001\,x_{10} - 0.19\,\text{Auto}
        \]

        \begin{center}
            \small
            \begin{tabular}{lrrrrr}
                \toprule
                \textbf{Variable} & \textbf{Coeff.} & \textbf{SE} & \textbf{t-stat} &
                    \textbf{p-value} \\
                \midrule
                Intercept   &  34.50 & 3.11 & 11.08 & $< 0.001$ \\
                $x_1$ (Displacement) & $-0.047$ & 0.026 & $-1.85$ & 0.076 \\
                $x_2$ (Horsepower)   & $+0.025$ & 0.042 & $+0.60$ & 0.552 \\
                $x_{10}$ (Weight)    & $-0.001$ & 0.002 & $-0.75$ & 0.460 \\
                Auto                 & $-0.193$ & 2.448 & $-0.08$ & 0.938 \\
                \midrule
                \multicolumn{4}{l}{$F_{\text{cal}} = 24.17$, $p < 0.001$;
                    \quad $R^2 = 0.782$} \\
                \bottomrule
            \end{tabular}
        \end{center}
    \end{frame}


    \begin{frame}{Multiple Regression --- Example: Interpretation}
        \begin{itemize}
            \item \textbf{Overall $F$-test:} $F_{\text{cal}} = 24.17 \gg F_{0.05,\,4,\,27}$;
                \textbf{reject} $H_0$ --- the model is statistically significant overall
            \item \textbf{Individual $t$-tests at $\alpha = 0.05$:}
            \begin{itemize}
                \item All four predictors have $p > 0.05$ --- individually,  none is significant at 5\% when the others are present
                \item This apparent contradiction (significant $F$, non-significant $t$'s) can occur due to \textbf{multicollinearity} among predictors
            \end{itemize}
            \item \textbf{Goodness of fit:} $R^2 = 0.782$ --- 78.2\% of the variation in mileage is explained by the model
            \item \textbf{Prediction} ($x_1=25$, $x_2=350$, $x_{10}=3800$, Auto):
                \[
                    \hat{y} = 34.5 - 0.05(25) + 0.03(350) - 0.001(3800) - 0.19(1)
                    \approx 39.76 \text{ mpg}
                \]
        \end{itemize}
    \end{frame}


    %------------------------------------------------------------
    \section{Indicator Variables}
    %------------------------------------------------------------

    \begin{frame}{Indicator (Dummy) Variables}
        \begin{itemize}
            \item Regression requires numerical inputs, but predictors may be \textbf{qualitative} (categorical), e.g.\ transmission type, tool type, sex
            \item A qualitative variable with $a$ categories is encoded using $(a - 1)$ \textbf{indicator (dummy) variables}, each taking values 0 or 1
            \item \textbf{Example} --- tool type (A, B, C):
                \[
                    x_B = \begin{cases} 1 & \text{if tool type B} \\ 0 & \text{otherwise} \end{cases},
                    \qquad
                    x_C = \begin{cases} 1 & \text{if tool type C} \\ 0 & \text{otherwise} \end{cases}
                \]
            \item Model: $y = \beta_0 + \beta_1 x_1 + \beta_2 x_B + \beta_3 x_C + \epsilon$
            \item Category A is the \textbf{reference group} ($x_B = x_C = 0$); $\beta_2$ and $\beta_3$ measure the \emph{difference} in mean $Y$ relative to category A
        \end{itemize}
    \end{frame}


    \begin{frame}{The Problem with $R^2$ in Multiple Regression}
        \begin{itemize}
            \item $R^2$ \textbf{never decreases} when a new predictor is added to the model --- even if that predictor is completely irrelevant
            \item Adding any variable, even random noise, will slightly increase (or at worst leave unchanged) $R^2$
            \item \textbf{Why?} Each new predictor uses up one more degree of freedom to ``fit'' by chance alone
            \item This means $R^2$ \textbf{rewards model complexity} and cannot be used to compare models with different numbers of predictors
            \item \textbf{Example:} adding 10 random noise variables to a model will still raise  $R^2$, even though the model has not improved We need a measure that \textbf{penalises} unnecessary predictors $\Rightarrow$ \textbf{Adjusted $R^2$}
        \end{itemize}
    \end{frame}


    \begin{frame}{Adjusted $R^2$}
        \textbf{Formula:}
        \[
            R^2_{\text{adj}}
            = 1 - \frac{\text{SSE}/(n - k - 1)}{\text{TSS}/(n - 1)}
            = 1 - \left(\frac{n-1}{n-k-1}\right)(1 - R^2)
        \]
        where $k$ is the number of predictors in the model.

        \vspace{0.5em}

        \begin{itemize}
            \item The penalty factor $\dfrac{n-1}{n-k-1} > 1$ increases as $k$ grows, \textbf{discounting} $R^2$ for model complexity
            \item Adding a predictor increases $R^2_{\text{adj}}$ \emph{only if} that predictor reduces MSE, i.e.\ if it genuinely improves fit beyond what is expected by chance
            \item $R^2_{\text{adj}} \leqslant R^2$; it can be negative for very poor models
            \item \textbf{Use $R^2_{\text{adj}}$} when comparing models with different numbers of predictors; use $R^2$ when reporting the fit of a single, already-selected model
        \end{itemize}
    \end{frame}


    %------------------------------------------------------------
    \section{Steps and Uses of Regression Analysis}
    %------------------------------------------------------------

    \begin{frame}{Steps of Regression Analysis}
        \begin{enumerate}
            \item \textbf{Hypothesize} a model of relationship (e.g.\ linear, multiple, logistic)
            \item \textbf{Estimate} the regression equation (LSE / OLS)
            \item \textbf{Test hypotheses} about the parameters ($F$-test for overall significance; $t$-tests for individual coefficients)
            \item Evaluate \textbf{goodness of fit} ($R^2$, residual plots)
            \item Use the model for \textbf{prediction}
        \end{enumerate}
    \end{frame}


    \begin{frame}{Uses of Regression Analysis}
        \begin{itemize}
            \item \textbf{Estimate} the average relationship between $Y$ and the predictors
            \item \textbf{Identify} which explanatory variables significantly affect $Y$, controlling for others
            \item \textbf{Predict} the value of $Y$ for given values of the predictors
            \item \textbf{Other regression types} (depending on $Y$):
            \begin{center}
                \small
                \begin{tabular}{ll}
                    \toprule
                    \textbf{Type of $Y$} & \textbf{Model} \\
                    \midrule
                    Continuous (normal) & Linear regression \\
                    Binary (0/1) & Logistic regression \\
                    Count & Poisson regression \\
                    Multiple correlated $Y$'s & Multivariate regression \\
                    \bottomrule
                \end{tabular}
            \end{center}
        \end{itemize}
    \end{frame}


    % TODO: add regression diagnostics


    \section*{Thank you.}

\end{document}